
\documentclass[11pt]{article}
\usepackage[english]{babel}
\usepackage[a4paper,margin=0.9in]{geometry}
\usepackage[T1]{fontenc}
\usepackage[utf8]{inputenc}
\usepackage{lmodern}
\usepackage{amsmath,amssymb,mathtools}
\usepackage{microtype}
\usepackage{fancyhdr}
\usepackage{array,enumitem}
\setlength{\parindent}{1.4em}
\setlength{\parskip}{0.15em}
\emergencystretch=8em
\setlength{\headheight}{14pt}
\pagestyle{fancy}
\fancyhf{}
\cfoot{\thepage}
\newcommand{\sect}[2]{\section*{\S\,#1. #2}}
\newcommand{\D}{\mathrm D}
\newcommand{\eps}{\varepsilon}
\lhead{Weber, Lehrbuch der Algebra I}
\rhead{English translation: \S\,73--\S\,75}
\begin{document}
\begin{center}
{\large Heinrich Weber}\\[0.35em]
{\Large Textbook of Algebra. Volume I.}\\[0.35em]
{\large Sixth section.}\\[0.25em]
{\large \S\,73--\S\,75}
\end{center}

\sect{73}{The Bezoutian}

When equation (4) of the preceding paragraph is expanded it has the form
\[
 y^n+P_1y^{n-1}+P_2y^{n-2}+\cdots+P_n=0,
\]
or, if $t_{n-1}$ is chosen so that $S(y)=0$,
\begin{equation}
 y^n+P_2y^{n-2}+\cdots+P_n=0.
\tag{1}
\end{equation}
Then
\[
 -2P_2=S(y^2)
\]
is a function of the second degree in the $t$'s and also in the $a$'s. This function is especially important for many applications and has received from Sylvester the name Bezoutian of the function $f(x)$, in honour of the French mathematician Bezout, who already in the preceding century gave the first correct developments of elimination.

By the formulas of the preceding paragraph this function can be computed rather simply. Beside the variables $t_{n-1},t_{n-2},\ldots,t_0$ we introduce a second independent system of variables $\tau_{n-1},\tau_{n-2},\ldots,\tau_0$ and put
\begin{equation}
\begin{aligned}
 y&=t_{n-1}f_0+t_{n-2}f_1+\cdots+t_0f_{n-1},\\
 z&=\tau_{n-1}f_0+\tau_{n-2}f_1+\cdots+\tau_0f_{n-1}.
\end{aligned}
\tag{2}
\end{equation}
Instead of forming $S(y^2)$ at once, we first compute $S(yz)$; $S(y^2)$ is then obtained by putting the $\tau$'s equal to the corresponding $t$'s.

Multiplying the formula
\begin{equation}
 yf_s=E_{0,s}f_0+E_{1,s}f_1+\cdots+E_{n-1,s}f_{n-1}
\tag{3}
\end{equation}
of the preceding paragraph by $\tau_{n-s-1}$ and summing with respect to $s$ from $0$ to $n-1$, one obtains
\begin{equation}
 yz=f_0\sum_{s=0}^{n-1}E_{0,s}\tau_{n-s-1}
   +f_1\sum_{s=0}^{n-1}E_{1,s}\tau_{n-s-1}
   +\cdots+f_{n-1}\sum_{s=0}^{n-1}E_{n-1,s}\tau_{n-s-1}.
\tag{4}
\end{equation}
Now by \S\,68, (6)
\[
 Sf_0=na_0,\qquad Sf_1=(n-1)a_1,\qquad \ldots,\qquad Sf_{n-1}=a_{n-1},
\]
and therefore
\[
\begin{aligned}
S(yz)=&\;na_0\sum_{s=0}^{n-1}E_{0,s}\tau_{n-s-1}
+(n-1)a_1\sum_{s=0}^{n-1}E_{1,s}\tau_{n-s-1}\\
&+\cdots+a_{n-1}\sum_{s=0}^{n-1}E_{n-1,s}\tau_{n-s-1}.
\end{aligned}
\]
In this expression the coefficient of $\tau_{n-1}$ is
\[
 na_0E_{0,0}+(n-1)a_1E_{1,0}+\cdots+a_{n-1}E_{n-1,0},
\]
and hence, by (3), is
\[
 a_0S(y).
\]
Since $S(z)$ begins with the term $na_0\tau_{n-1}$, the difference
\begin{equation}
 S(yz)-\frac1nS(y)S(z)
\tag{5}
\end{equation}
contains no term multiplied by $\tau_{n-1}$. Since this expression is also symmetric in $t$ and $\tau$, it also contains no $t_{n-1}$; in forming it we may therefore simply put $t_{n-1}$ and $\tau_{n-1}$ equal to zero.

Put
\[
 S(y)=na_0t_{n-1}+\varphi(t),
\]
where
\begin{equation}
 \varphi(t)=(n-1)a_1t_{n-2}+(n-2)a_2t_{n-3}+\cdots+a_{n-1}t_0.
\tag{6}
\end{equation}
Then
\begin{equation}
\begin{aligned}
S(yz)-\frac1nS(y)S(z)=&\;na_0\sum_{s=1}^{n-1}E_{0,s}\tau_{n-s-1}
+(n-1)a_1\sum_{s=1}^{n-1}E_{1,s}\tau_{n-s-1}\\
&+\cdots+a_{n-1}\sum_{s=1}^{n-1}E_{n-1,s}\tau_{n-s-1}
-\frac1n\varphi(t)\varphi(\tau).
\end{aligned}
\tag{7}
\end{equation}
Here it is to be observed that $t_{n-1}=0$ is to be assumed; thus
\[
 E_{s,s}=a_1t_{n-2}+\cdots+a_st_{n-s-1}
\]
is to be put, while the remaining $E$'s are determined by the formulas (9) of the preceding paragraph.

The expression on the right of (7) is a bilinear function of $t$ and $\tau$. Let its expanded form be
\[
 \sum_{h=0}^{n-2}\sum_{k=0}^{n-2}B_{h,k}t_h\tau_k,
\]
with $B_{i,k}=B_{k,i}$. The Bezoutian is then obtained by putting $\tau=t$, namely
\begin{equation}
 B=\sum_{h=0}^{n-2}\sum_{k=0}^{n-2}B_{h,k}t_ht_k,
\tag{8}
\end{equation}
where the coefficients $B_{i,k}$ are quadratic functions of the $a$'s. From (7) one obtains
\begin{equation}
 S(yz)-\frac1nS(y)S(z)=\sum_{h=0}^{n-2}\sum_{k=0}^{n-2}B_{h,k}t_h\tau_k,
\tag{9}
\end{equation}
\begin{equation}
 S(y^2)=\frac1n[S(y)]^2+B.
\tag{10}
\end{equation}
By formula (8) of \S\,68,
\[
\begin{aligned}
 y-\frac1nS(y)&=t_{n-2}F_0+t_{n-3}F_1+\cdots+t_0F_{n-2},\\
 z-\frac1nS(z)&=\tau_{n-2}F_0+\tau_{n-3}F_1+\cdots+\tau_0F_{n-2}.
\end{aligned}
\]
Multiplying these two expressions gives
\[
 yz-\frac1n[yS(z)+zS(y)]+\frac1{n^2}S(y)S(z)
 =\sum_{h,k=0}^{n-2}\tau_ht_kF_{n-h-2}F_{n-k-2}.
\]
Summing this formula over all roots $x$, that is, taking the sum denoted by $S$, gives
\[
 S(yz)-\frac1nS(y)S(z)=\sum_{h,k=0}^{n-2}\tau_ht_kS(F_{n-h-2}F_{n-k-2}).
\]
Comparison with (9) yields
\begin{equation}
 B_{h,k}=S(F_{n-h-2}F_{n-k-2}).
\tag{11}
\end{equation}
Let $\Delta$ denote the determinant of the function $B$:
\[
 \Delta=\sum \pm B_{0,0}B_{1,1}\cdots B_{n-2,n-2}.
\]
By the multiplication theorem for determinants, and using
\[
 SF_0=0,
 \qquad SF_1=0,
 \qquad\ldots,
 \qquad SF_{n-2}=0,
\]
if
\[
 \Phi=\begin{vmatrix}
 1&F_0(x_1)&F_1(x_1)&\cdots&F_{n-2}(x_1)\\
 1&F_0(x_2)&F_1(x_2)&\cdots&F_{n-2}(x_2)\\
 \vdots&\vdots&\vdots&&\vdots\\
 1&F_0(x_n)&F_1(x_n)&\cdots&F_{n-2}(x_n)
 \end{vmatrix},
\]
then (11) gives
\begin{equation}
 \Phi^2=n\Delta.
\tag{12}
\end{equation}
Now using the expressions (9) of \S\,68 for $F_0,F_1,\ldots,F_{n-2}$, and the theorem that in a determinant one may add to one column another column multiplied by an arbitrary factor (\S\,22, VII), $\Phi$ is transformed into
\[
 a_0^{n-1}
 \begin{vmatrix}
 1&x_1&x_1^2&\cdots&x_1^{n-1}\\
 1&x_2&x_2^2&\cdots&x_2^{n-1}\\
 \vdots&\vdots&\vdots&&\vdots\\
 1&x_n&x_n^2&\cdots&x_n^{n-1}
 \end{vmatrix},
\]
the square of which is, by \S\,46, the discriminant of $f(x)$. Thus (12) gives the important theorem:

The determinant of the Bezoutian of $f(x)$ is the $n$th part of the discriminant of $f(x)$.

By these formulas the computation of the Bezoutian is no longer difficult and is quite simple. For $n=3$ one obtains the result already to be inferred from the formulas of \S\,71,
\begin{equation}
 B=-\frac23H(t_1,t_0).
\tag{13}
\end{equation}

To illustrate the preceding discussion, we write out the complete formula system for the computation of the Bezoutian of the biquadratic form, leaving to the reader the execution of the only slightly longer calculation for the form of the fifth degree. For the form of the fourth degree we have, from (7) and \S\,72, (9),
\[
 f(x)=a_0x^4+a_1x^3+a_2x^2+a_3x+a_4,
\]
\[
\begin{array}{rl|l}
E_{0,1}={}&-a_2t_2-a_3t_1-a_4t_0&4a_0t_2\\
E_{1,1}={}&a_1t_2&3a_1t_2\\
E_{2,1}={}&a_0t_2+a_1t_1&2a_2t_2\\
E_{3,1}={}&a_0t_1+a_1t_0&a_3t_2\\ \hline
E_{0,2}={}&-a_3t_2-a_4t_1&4a_0t_1\\
E_{1,2}={}&-a_3t_1-a_4t_0&3a_1t_1\\
E_{2,2}={}&a_1t_2+a_2t_1&2a_2t_1\\
E_{3,2}={}&a_0t_2+a_1t_1+a_2t_0&a_3t_1\\ \hline
E_{0,3}={}&-a_4t_2&4a_0t_0\\
E_{1,3}={}&-a_4t_1&3a_1t_0\\
E_{2,3}={}&-a_4t_0&2a_2t_0\\
E_{3,3}={}&a_1t_2+a_2t_1+a_3t_0&a_3t_0.
\end{array}
\]
These expressions are to be multiplied by the factors written to their right and added; then one must add
\[
 -\frac14\varphi(t)\varphi(\tau)
 =-\frac14(3a_1t_2+2a_2t_1+a_3t_0)(3a_1\tau_2+2a_2\tau_1+a_3\tau_0),
\]
and collect the coefficients of $\tau_ht_k$. This gives
\begin{equation}
\begin{aligned}
B_{2,2}&=-\frac14(8a_0a_2-3a_1^2),&
B_{0,1}&=-\frac12(6a_1a_4-a_2a_3),\\
B_{1,1}&=-4a_0a_4-2a_1a_3+a_2^2,&
B_{1,2}&=-\frac12(6a_0a_3-a_1a_2),\\
B_{0,0}&=-\frac14(8a_2a_4-3a_3^2),&
B_{0,2}&=-\frac14(16a_0a_4-a_1a_3).
\end{aligned}
\tag{14}
\end{equation}
In the same way one may proceed with a form of the fifth degree, and obtains for the coefficients of the Bezoutian
\[
\begin{aligned}
5B_{0,0}&=4a_4^2-10a_3a_5,\\
5B_{1,1}&=6a_3^2-10a_2a_4-20a_1a_5,\\
5B_{2,2}&=6a_2^2-10a_1a_3-20a_0a_4,\\
5B_{3,3}&=4a_1^2-10a_0a_2,
\end{aligned}\qquad
\begin{aligned}
5B_{0,1}&=3a_3a_4-15a_2a_5,\\
5B_{0,2}&=2a_2a_4-20a_1a_5,\\
5B_{0,3}&=a_1a_4-25a_0a_5,\\
5B_{1,2}&=4a_2a_3-15a_1a_4-25a_0a_5,\\
5B_{1,3}&=2a_1a_3-20a_0a_4,\\
5B_{2,3}&=3a_1a_2-15a_0a_3.
\end{aligned}
\]

\sect{74}{Transformation of the equation of the fifth degree}

These developments will now be applied in order to carry out the transformation of the equation of the fifth degree as it was sketched in \S\,54.

In infinitely many ways a system of values $t_0,t_1,t_2,t_3$ can be so determined that the Bezoutian $B$ vanishes. In general this requires the solution of a quadratic equation. For example, we may put
\[
 t_2=0,
 \qquad t_3=0
\]
and determine the ratio $t_0:t_1$ from the quadratic equation
\[
 B_{0,0}t_0^2+2B_{0,1}t_0t_1+B_{1,1}t_1^2=0.
\]
These values of $t$ then contain the square root
\[
 \sqrt{B_{0,1}^2-B_{0,0}B_{1,1}}.
\]
Instead one can make other determinations and obtains other square roots. In the choice of this square root there is something arbitrary and indeterminate.

Among these various possible determinations we choose one and thereby transform the given equation of the fifth degree into a principal equation, that is, into one in which the third and fourth powers of the unknown do not occur; or we suppose, following Klein's procedure, that the equation to be transformed is already from the outset a principal equation
\begin{equation}
 f(x)=a_0x^5+a_3x^2+a_4x+a_5=0.
\tag{1}
\end{equation}
We must, however, expressly emphasize that by \S\,70, (8), under this preliminary Tschirnhausen transformation the discriminant of the given equation is changed only by a factor $\Theta^2$, where $\Theta$ is rational in the applied $t$'s and hence also in the square root used in determining them.

Under the supposition (1), the coefficients of the Bezoutian are, by the formulas of the preceding paragraph,
\begin{equation}
\begin{array}{ll}
5B_{0,0}=4a_4^2-10a_3a_5, & 5B_{0,1}=3a_3a_4,\\
5B_{1,1}=6a_3^2, & 5B_{0,2}=0,\\
5B_{2,2}=-20a_0a_4, & 5B_{0,3}=-25a_0a_5,\\
5B_{3,3}=0, & 5B_{1,2}=-25a_0a_5,\\
5B_{1,3}=-20a_0a_4, & 5B_{2,3}=-15a_0a_3.
\end{array}
\tag{2}
\end{equation}
Computing the determinant from these entries, one obtains without difficulty, by \S\,73, (12), the discriminant of the principal equation of the fifth degree:
\begin{equation}
\begin{aligned}
D=a_0^5(&5^5a_0^3a_5^4+2^8a_0a_4^5+2250a_0a_3^2a_4a_5^2\\
&-1600a_0a_3a_4^3a_5+108a_3^5a_5-27a_3^4a_4^2).
\end{aligned}
\tag{3}
\end{equation}
Since here $B_{3,3}=0$, we already have a system of values of the $t$ for which the Bezoutian vanishes, namely
\[
 t_0=0,
 \qquad t_1=0,
 \qquad t_2=0.
\]
We now take our Tschirnhausen substitution in the form [\S\,68, (12)]
\begin{equation}
 y=t_3F_0+t_2(\alpha_2F_1+\alpha_1F_2+\alpha_0F_3),
\tag{4}
\end{equation}
and choose $\alpha_0,\alpha_1,\alpha_2$ so that the equation of the fifth degree for $y$ becomes, independently of $t_3,t_2$, a principal equation, that is, so that $S(y^2)$ vanishes identically for all $t_2,t_3$.

Since $S(F_0^2)=B_{3,3}$ vanishes here, the requirement is satisfied, by \S\,73, (11), if $\alpha_0,\alpha_1,\alpha_2$ are chosen so that
\begin{equation}
\begin{aligned}
&\alpha_0^2B_{0,0}+\alpha_1^2B_{1,1}+\alpha_2^2B_{2,2}
+2\alpha_0\alpha_1B_{0,1}+2\alpha_0\alpha_2B_{0,2}+2\alpha_1\alpha_2B_{1,2}=0,
\end{aligned}
\tag{5}
\end{equation}
\begin{equation}
 \alpha_0B_{0,3}+\alpha_1B_{1,3}+\alpha_2B_{2,3}=0.
\tag{6}
\end{equation}
By eliminating one of the three quantities $\alpha$, these equations lead to a quadratic equation for the ratio of the other two. In itself there is no need to exclude any special case, not even the case in which $B_{0,3},B_{1,3},B_{2,3}$ all vanish; one would then merely have to choose any solution of equation (5) for the $\alpha$'s.

In order not to be prolix, however, we assume that $B_{0,3}$ is non-zero. The cases thereby excluded, in which a root of the given equation becomes zero or infinite, are of no real interest here, since they reduce to an equation of lower degree. The treatment remains the same if one assumes instead that another of $B_{0,3},B_{1,3},B_{2,3}$ is non-zero.

To treat equations (5), (6) symmetrically, and especially to recognize which square root is required for their solution, we proceed as follows. Put, for abbreviation,
\begin{equation}
\begin{aligned}
B_0&=\alpha_0B_{0,0}+\alpha_1B_{0,1}+\alpha_2B_{0,2},\\
B_1&=\alpha_0B_{1,0}+\alpha_1B_{1,1}+\alpha_2B_{1,2},\\
B_2&=\alpha_0B_{2,0}+\alpha_1B_{2,1}+\alpha_2B_{2,2}.
\end{aligned}
\tag{7}
\end{equation}
Then equations (5) and (6) can be written as
\begin{equation}
\begin{aligned}
\alpha_0B_0+\alpha_1B_1+\alpha_2B_2&=0,\\
\alpha_0B_{0,3}+\alpha_1B_{1,3}+\alpha_2B_{2,3}&=0.
\end{aligned}
\tag{8}
\end{equation}
We multiply the second equation by an undetermined factor $\alpha_3$ and add it to the first, obtaining
\begin{equation}
 \alpha_0(B_0+\alpha_3B_{0,3})+
 \alpha_1(B_1+\alpha_3B_{1,3})+
 \alpha_2(B_2+\alpha_3B_{2,3})=0.
\tag{9}
\end{equation}
This equation is satisfied if we put
\begin{equation}
 B_0+\alpha_3B_{0,3}=0,
 \qquad B_1+\alpha_3B_{1,3}=\sigma\alpha_2,
 \qquad B_2+\alpha_3B_{2,3}=-\sigma\alpha_1.
\tag{10}
\end{equation}
From these three equations, together with (6), the unknowns $\alpha_0,\alpha_1,\alpha_2,\alpha_3,\sigma$ are to be determined. Written out, the equations are
\begin{equation}
\begin{aligned}
\alpha_0B_{0,0}+\alpha_1B_{0,1}+\alpha_2B_{0,2}+\alpha_3B_{0,3}&=0,\\
\alpha_0B_{1,0}+\alpha_1B_{1,1}+\alpha_2(B_{1,2}-\sigma)+\alpha_3B_{1,3}&=0,\\
\alpha_0B_{2,0}+\alpha_1(B_{2,1}+\sigma)+\alpha_2B_{2,2}+\alpha_3B_{2,3}&=0,\\
\alpha_0B_{3,0}+\alpha_1B_{3,1}+\alpha_2B_{3,2}&=0.
\end{aligned}
\tag{11}
\end{equation}
Eliminating $\alpha_0,\alpha_1,\alpha_2,\alpha_3$ gives a quadratic equation for $\sigma$; after solving it the ratios $\alpha_0:\alpha_1:\alpha_2:\alpha_3$ can be determined from linear equations.

The quadratic equation for $\sigma$ is, in determinant form,
\[
\begin{vmatrix}
 B_{0,0}&B_{0,1}&B_{0,2}&B_{0,3}\\
 B_{1,0}&B_{1,1}&B_{1,2}-\sigma&B_{1,3}\\
 B_{2,0}&B_{2,1}+\sigma&B_{2,2}&B_{2,3}\\
 B_{3,0}&B_{3,1}&B_{3,2}&0
\end{vmatrix}=0.
\]
Since the left hand side is unchanged when $\sigma$ is replaced by $-\sigma$, it is a pure quadratic equation, and by \S\,73, (12) it gives
\[
 B_{0,3}^2\sigma^2-5D=0,
\]
where $D$ is the discriminant of the given equation. Hence by (2) we obtain
\begin{equation}
 5a_0a_5\sigma=\sqrt{5D},
\tag{12}
\end{equation}
with either choice of sign.

\sect{75}{Normal form of the equation of the fifth degree}

As was remarked earlier, the principal aim of these considerations is to produce a normal form of the equation of the fifth degree which depends on only one undetermined coefficient, one parameter. One such normal form is the Bring-Jerrard form. To obtain it, according to (4) of the preceding paragraph one must put
\begin{equation}
 y=t_3F_0+t_2(\alpha_2F_1+\alpha_1F_2+\alpha_0F_3),
\tag{1}
\end{equation}
so that identically
\[
 S(y)=0,
 \qquad S(y^2)=0,
\]
and then determine the ratio $t_3:t_2$ from the cubic equation
\[
 S(y^3)=0.
\]
This cubic equation can indeed be formed, although its expression is long. The formulas required for it were computed by Cayley.

It is the result of a remarkable investigation of Gordan that another normal form, the Brioschi normal form, can be obtained without any new irrationality. We shall derive this result at the close of the discussion of the Tschirnhausen transformation.

We retain the assumptions of the preceding paragraph and put, for the moment,
\begin{equation}
 u=F_0,
 \qquad
 v=\alpha_2F_1+\alpha_1F_2+\alpha_0F_3,
\tag{2}
\end{equation}
where $\alpha_0,\alpha_1,\alpha_2$ are to have the values determined in the preceding paragraph.

By formula \S\,73, (4), the three functions
\[
 u^2,
 \qquad uv,
 \qquad v^2
\]
can be represented linearly and homogeneously in terms of $f_0,f_1,f_2,f_3,f_4$; and since
\begin{equation}
 S(u^2)=0,
 \qquad S(uv)=0,
 \qquad S(v^2)=0,
\tag{3}
\end{equation}
these expressions are also linear and homogeneous in $F_0,F_1,F_2,F_3$.

The calculation is easily carried out from the formulas of \S\S\,72, 73, but will not be pursued here, since only the underlying idea is needed. We record only that the coefficients in the expressions for $u^2,uv,v^2$ are linear in the coefficients of the original equation of the fifth degree and quadratic in $\alpha_0,\alpha_1,\alpha_2$. If we now eliminate $F_0,F_1,F_2,F_3$ from these expressions using (2), we obtain a relation of the form
\begin{equation}
 pu^2+2quv+rv^2=au+bv,
\tag{4}
\end{equation}
where $p,q,r,a,b$ depend rationally on the coefficients $a_i$ and $\alpha_i$, and in any case do not all vanish simultaneously.

For the moment put
\[
 pu^2+2quv+rv^2=\varphi(u,v).
\]
Then by \S\,57,
\[
 4\varphi(u',v')\varphi(u,v)-[u'\varphi'(u)+v'\varphi'(v)]^2
\]
is the square of a linear function of $u,v$, so that $\varphi(u,v)$ is thereby decomposed into the sum of two squares. Choosing $u'=b$, $v'=-a$, one obtains
\[
 \varphi(b,-a)\varphi(u,v)=[b(pu+qv)-a(qu+rv)]^2+(pr-q^2)(au+bv)^2.
\]
Thus, with the abbreviations
\begin{equation}
\begin{aligned}
pb^2-2qab+ra^2&=m,\\
bp-aq&=a',\\
bq-ar&=b',\\
q^2-pr&=c,
\end{aligned}
\tag{5}
\end{equation}
we get
\begin{equation}
 m(pu^2+2quv+rv^2)=(a'u+b'v)^2-c(au+bv)^2.
\tag{6}
\end{equation}
By (4) this equation may also be written as
\begin{equation}
 m(au+bv)=(a'u+b'v)^2-c(au+bv)^2.
\tag{7}
\end{equation}
Now put
\begin{equation}
 y=\frac{a'u+b'v}{au+bv},
\tag{8}
\end{equation}
and form the equation of the fifth degree whose roots are the five values of $y$:
\begin{equation}
 y^5+c_1y^4+c_2y^3+c_3y^2+c_4y+c_5=0.
\tag{9}
\end{equation}
The coefficients $c$ can then be further determined as follows.

From (8) one obtains
\[
 y+\sqrt c=\frac{a'u+b'v+\sqrt c(au+bv)}{au+bv},
\]
and hence, by (7),
\[
 y+\sqrt c=\frac{m}{a'u+b'v-\sqrt c(au+bv)},
\qquad
\frac{m}{y+\sqrt c}=a'u+b'v-\sqrt c(au+bv).
\]
It follows from (3) that
\[
 S\left(\frac1{y+\sqrt c}\right)=0,
 \qquad
 S\left(\frac1{(y+\sqrt c)^2}\right)=0.
\]
If from (9) we derive the equation whose roots are the values
\[
 \xi=-\frac1{y+\sqrt c},
\]
that is, if we put
\[
 y=\frac{1-\xi\sqrt c}{\xi},
\]
then the resulting equation in $\xi$ must have vanishing coefficients of $\xi^4$ and $\xi^3$, whichever sign is assigned to the square root $\sqrt c$. This yields four equations among the coefficients $c_\nu$ of (9).

The equation in $\xi$ is
\[
\begin{aligned}
&(1-\xi\sqrt c)^5+c_1\xi(1-\xi\sqrt c)^4+c_2\xi^2(1-\xi\sqrt c)^3\\
&\qquad +c_3\xi^3(1-\xi\sqrt c)^2+c_4\xi^4(1-\xi\sqrt c)+c_5\xi^5=0.
\end{aligned}
\]
Setting the coefficients of $\xi^4$ and $\xi^3$ equal to zero gives
\[
\begin{aligned}
5\sqrt c^{\,4}-4c_1\sqrt c^{\,3}+3c_2\sqrt c^{\,2}-2c_3\sqrt c+c_4&=0,\\
-10\sqrt c^{\,3}+6c_1\sqrt c^{\,2}-3c_2\sqrt c+c_3&=0.
\end{aligned}
\]
Because of the two possible signs of $\sqrt c$, these equations split into the four equations
\[
\begin{aligned}
5c^2+3c_2c+c_4&=0,\\
4c_1c+2c_3&=0,\\
10c+3c_2&=0,\\
6c_1c+c_3&=0.
\end{aligned}
\]
The second and fourth give
\[
 c_1=0,
 \qquad c_3=0,
\]
and then the third and first give
\[
 c_2=-\frac{10}{3}c,
 \qquad c_4=-5c^2.
\]
Thus the equation for $y$ assumes the form
\begin{equation}
 y^5-\frac{10}{3}cy^3+5c^2y+c_5=0.
\tag{10}
\end{equation}
This equation still depends on the two parameters $c_5$ and $c$. It can be reduced to an equation with one parameter by the substitution
\begin{equation}
 y=\sqrt{\frac c3}\,z,
 \qquad
 \gamma=c_5\sqrt{\frac{3^5}{c^5}},
\tag{11}
\end{equation}
which gives the simple and elegant form
\begin{equation}
 z^5-10z^3+45z+\gamma=0.
\tag{12}
\end{equation}
This is the Brioschi normal form. The substitution (11), however, has the disadvantage that it still contains a square root, whereas in formulas (8) and (10) only the two square roots discussed in the preceding paragraph occur. But, by a remark of Klein, one can also pass rationally from (10) to a normal form containing only one parameter.

For example, putting
\begin{equation}
 y=\frac{3c_5}{c^2}\frac1z,
 \qquad
 \gamma=\frac{9c_5^2}{c^5},
\tag{13}
\end{equation}
one obtains from (10) the equation
\begin{equation}
 z^5+15z^4-10\gamma z^2+3\gamma^2=0.
\tag{14}
\end{equation}
In this transformation it has been tacitly assumed that the quantities denoted in (5) by $m$ and $c$ are not zero.

If $m=0$, then in formula (4) the expression $pu^2+2quv+rv^2$ is divisible by $au+bv$; hence at least for some of the roots the equation
\[
 au+bv=0
\]
must hold, that is, an equation of degree not higher than the fourth.

If $c=0$, so that $pu^2+2quv+rv^2$ is the square of a linear function, then (7) shows that, on putting
\[
 y=a'u+b'v,
\]
one has
\[
 S(y)=0,
 \qquad S(y^2)=0,
 \qquad S(y^3)=0,
\]
and therefore the Bring-Jerrard form can be produced rationally.

We may finally summarize the results of these considerations as follows:

The principal equation of the fifth degree can be transformed into a normal form with one parameter by a transformation whose only irrationality is the square root of the discriminant.

\end{document}
